\section{Problem Formulation}
\label{sec:problem}

Structured evidence pipelines usually change two things at once: \emph{which facts} reach the answerer, and \emph{how} those facts are arranged and annotated. A transformation can lose a fact (an extractor omits it, a budget cuts it), and it can add information that is not a fact of the source text (an order, a grouping into steps, a relation between two facts). We measure the two separately.

\paragraph{Facts and representations.}
For a question $q$ with gold answers $A_q$ we start from a fixed set of source passages; in our controlled experiments these are the dataset's gold supporting passages. Segmenting them into sentences gives the fact set $F_q=\{f_1,\dots,f_n\}$, where each fact keeps the identifier of its source span. Because facts are sentences of the gold passages, not model-generated paraphrases, $F_q$ contains no extraction error. A sentence is not an atomic proposition, however: some sentences need their passage to be read (\emph{It stars \dots}), so segmentation can make a fact harder to interpret without removing it. Model-based proposition extraction enters only through the retrieved-evidence systems of Experiment~0, and we never mix extracted claims into an information-controlled comparison. A representation $\rho$ maps the facts, optionally with structure $S_q$ derived from annotations or predicted by a model, to an evidence block
\begin{equation}
E_q^{\rho}=\mathrm{render}_\rho\big(\pi_\rho(F_q),\,S_q\big),
\end{equation}
where $\pi_\rho$ orders and groups the facts and $\mathrm{render}_\rho$ serialises them, possibly with a list of relations. We write $\mathrm{facts}(E)$ for the facts that appear in a block, and $\mathrm{added}(\rho)$ for the information that the block carries beyond them. We distinguish three kinds of added information that structured pipelines usually change together: the \emph{order} $\pi_\rho$ of the items; explicit \emph{relation lines} $L_\rho\subseteq F_q\times F_q\times\mathcal{T}$, each asserting a typed relation (dependency, support, contradiction, duplication, supersession) between two items; and the \emph{provenance} of those relations, derived deterministically from dataset annotations or predicted by a model. Our controls change one of these at a time, and include a neutral label that asserts no relation type while keeping the same pairs and number of lines. A fixed answerer $M$ with one prompt template $\tau$, identical for every representation, produces $\hat a=M(\tau(q,E_q^{\rho}))$ with greedy decoding. The template allows the answer \emph{insufficient evidence}, so each outcome is one of three states: correct (C), answered but wrong (W), or not answered (A, an abstention or an unparseable reply). Exact match counts A as incorrect; we report the transitions between states separately, because a representation can change accuracy by changing the willingness to answer.

\paragraph{Protocol A: information-controlled.}
Protocol A compares representations whose fact sets are identical, $\mathrm{facts}(E_q^{\rho})=F_q$ for every $\rho$, which we verify per question by comparing identifiers and text digests. Under this constraint an accuracy difference cannot come from a missing fact; it can still come from $\mathrm{added}(\rho)$, which we record per record, and from formatting tokens. We record the content tokens $c_q=\sum_j |\mathrm{text}_j|$ and the overhead $o_q^{\rho}=|E_q^{\rho}|-c_q$ separately.

\paragraph{Protocol B: budget-controlled.}
Under a budget $B_q=\beta\, c_q$ we rank the facts once, with BM25 against the question, independently of the representation and without gold labels, and each condition keeps the longest prefix of the ranking whose rendering fits,
\begin{equation}
k_\rho=\max\{k:\ |\mathrm{render}_\rho(\pi_\rho(F_q^{\le k}),S_q)|\le B_q\}.
\end{equation}
Under the \emph{per-representation} policy each representation keeps its own $k_\rho$, so verbose formats keep fewer facts; under the \emph{shared} policy all keep $\min_\rho k_\rho$ and show the same facts. Facts are never cut mid-sentence. For every dropped fact we record whether it is a gold supporting fact and whether it contains the answer string; this is our measure of information loss.

\paragraph{Perturbation validity and objective.}
An intervention $\delta$ maps an evidence block $E$ to a corrupted block $\delta(E)$. It is \emph{label-preserving} if $A_q$ remains correct given $\delta(E)$. Reordering facts or relabelling a relation is label-preserving; deleting a required fact is not. Accuracy changes are estimated only on label-preserving records that do not introduce the answer through a distractor. For a contrast between $\rho$ and $\rho'$ we estimate the paired difference in question-level exact match over the questions $Q$ available in both conditions, with unanswered questions scored as incorrect:
\begin{equation}
\Delta_{\rho'\!-\!\rho}=\tfrac{1}{|Q|}\textstyle\sum_{q\in Q}\big(\mathrm{EM}(q,\rho')-\mathrm{EM}(q,\rho)\big).
\end{equation}
